\justifying
\NSFCBodyText

截至本轮材料核对时，申请人主持一项自治区高校基本科研业务费科研项目（培育类），题目为“融合多源数据大模型驱动的光伏发电功率预测方法研究”。该项目批准号为XJEDU2026P005，资助金额5万元（直接经费3.9万元、间接经费1.1万元），执行期2026年1月至12月，负责总体方案、模型研究与成果组织。本申请拟研究可用气象信息约束下的概率预测、天气转换和爬坡事件，不将仅更换场站或题目作为独立任务的证明；两项目任务、数据和成果边界须按研究基础所列F-XJ-05核验，并分别记录支出与成果归属。

申请人参与国家自然科学基金地区项目“数据驱动下基于互补型生成对抗网络与神经统计模型的公路沙害智能预警研究”，批准号62466056，资助金额32万元，执行期2025年1月至2028年12月，基金委登记单位为新疆大学，本人负责数据收集整理。该项目对象为公路沙害，本申请对象为光伏功率变化；数据质量控制经验可提供方法参考，但不将公路沙害资料视为已获授权的光伏场站数据，也不共用未获许可的成果。

上述内容依据已确认的共用申请人事实。拟参与人员及其相关在研项目须随本项目团队落实一并核对，不能以其他申报中的人员名单推定。本基金对未结题基本科研业务费等项目另有限制，研究内容有差异不等于自动满足申报资格；截至截止日的状态按F-XJ-05与正式审核材料处理。
